% !TEX root = ../thesis.tex
% \renewcommand{\labelenumii}{\arabic{enumi}.\arabic{enumii}}
% \renewcommand{\labelenumiii}{\arabic{enumi}.\arabic{enumii}.\arabic{enumiii}}
% \renewcommand{\labelenumiv}{\arabic{enumi}.\arabic{enumii}.\arabic{enumiii}.\arabic{enumiv}}

\mychapter{Introduction}{Introduction}{}
\label{chap:introduction}

\mysection{Background and Motivation}{Background and Motivation}

Accurately measuring human motion is a foundational challenge in sports science and rehabilitation. Specialists, such as strength and conditioning (SNC) coaches, traditionally rely on visual observation or 2D video data to assess technique and mitigate injury risks \cite{needham2021accuracy, chariar2023ai, forte2023biomechanics}.
Performing an exercise with poor form could contribute to biomechanical overload and an increased risk of overuse injuries, joint degeneration, and other conditions \cite{chariar2023ai}. \citet{forte2023biomechanics} state that proper exercise biomechanics involves the correct form and technique to ensure the joints, muscles and bones are moving safely. 
Analysing the biomechanics, movement patterns and movement metrics such as angles and alignment of joints can help specialists track and potentially mitigate these injuries \cite{forte2023biomechanics}.

With the advancement of technology, motion capture systems have been designed and tailored to assist in studying the biomechanics and kinematics of human motion \cite{colyer2018review, mundermann2006evolution}.
Optical motion capture systems can be divided into two broad categories namely marker-based and markerless motion capture.

Marker-based pose estimation systems rely on the placement of reflective markers on the human body and specialised hardware to capture the human
motion via infrared technology. Systems such as Vicon (Vicon Nexus, Vicon Motion Systems Ltd, Oxford, UK) are considered to be the industry gold standard for biomechanical applications \cite{matsuda2025reliability, hansen2024validation, song2022evolution}. However, marker-based systems have a number of drawbacks. These systems are expensive, restricted to laboratory environments, and require trained professionals knowledgeable in marker placement, system setup and calibration \cite{mercadal2024exercise,kuster2016accuracy}. Furthermore, the placement of markers can be invasive for the participants, requiring form-fitting clothing and exposed skin. These markers might affect the natural movement of the participant and their placement does not directly correspond to true anatomical joint centres as markers are used in conjunction with one another to infer joint centres \cite{needham2021accuracy}.   
These drawbacks have driven the exploration of markerless motion capture systems. 

Markerless pose estimation systems make use of optical capture cameras, human body models, computer vision and deep learning algorithms to capture and detect human motion \cite{colyer2018review,federolf2025validation}. They offer an attractive, non-invasive alternative to marker-based systems \cite{federolf2025validation}. Whilst markerless systems are not yet widely adopted, systems such as Theia3D (Theia Markerless Inc., Kingston, ON, Canada), have demonstrated accuracy comparable to Vicon for lower limb analysis, though further research is required to validate full-body biomechanics \cite{federolf2025validation, yang2025comparison, wren2023comparison,scataglini2024accuracy}.

While multiview markerless systems offer a promising alternative to marker-based systems their success is largely dependent on the quality, orientation and position of the cameras utilised \cite{rahimian2016optimal,jorgensen2018simulation, yang2025evaluation}.    
Literature has shown that cameras play a vital role in the success of pose estimation algorithms yet the full extent to which different configurations affect performance is not sufficiently understood \cite{kwak2025assessing}.

Finding the optimal camera configuration within the real world is impractical and time-consuming due to the configuration space being continuous. There are many variables to choose from that could influence a camera configuration's performance such as a camera's pan, tilt and location \cite{zhao2013approximate}. Physically testing configurations requires repeated setup and checkerboard calibration of each configuration being examined.
Given the infeasibility and difficulty of testing every physical multi-camera configuration in the real-world this physical optimisation problem motivates the use of simulation. Existing literature by \citet{zhao2013approximate}, \citet{rahimian2016optimal} and \citet{jorgensen2018simulation} has sought to find optimal camera placement for various contexts such as industrial pose estimation through simulation-based approaches, however they were outside the context of markerless human pose estimation.

It has also been highlighted that the choice of pose estimation algorithm could influence the choice of camera configuration \cite{samani2026optimizing}. 
A comparison of human pose estimation algorithms has not been performed on multi-view markerless camera configurations. It has been cited as a potential limitation and area for future work. This presents a gap in exploring how choice of human pose estimation algorithm influences system performance and potentially camera configuration \cite{samani2026optimizing,kim2026comparison}.

\citet{rahimian2016optimal} introduce and compare two methods for camera placement and utilise an optimisation algorithm to evaluate the two methods. 
\citet{jorgensen2018simulation} used a simulation-based approach but avoided photo-realistic rendering as it was complex and computationally expensive.
With the technological advancements of GPUs there has been a shift towards synthetic data generation as a means to augment and create datasets. Synthetic datasets have become increasingly popular within the Computer Vision community with early pioneering works like SURREAL by \citet{varol2017learning} and AGORA by \citet{patel2021agora} to highly detailed, photo-realistic, state-of-the-art (SOTA) datasets such as BEDLAM by \citet{black2023bedlam,tesch2025bedlam2} and SynBody by \citet{yang2023synbody}.

Simulation-based approaches utilising synthetic data have numerous benefits over real-world data. Primarily they provide known 3D human parameters and camera parameters such as the intrinsic and extrinsic matrices without the need for calibration \cite{black2023bedlam,jorgensen2018simulation}, repeatability and diverse subjects and environments \cite{black2023bedlam}. Using digital content creator (DCC) tools, researchers have access to exact locations of 3D joints extracted directly from the animated human's armature in contrast to the real-world inferred ground truth joint centres \cite{needham2021accuracy,patel2021agora}. Furthermore, using statistical parametric human body models such as the SMPL-X body model \cite{SMPL-X:2019}, it is feasible to alter subject phenotypes (height, weight, skin tone) allowing for highly diverse, repeatable datasets that mitigate demographic biases \cite{patel2021agora,black2023bedlam,yang2023synbody,bazavan2021hspace}.

However, whilst the simulation-based approach eliminates the physical constraints of real world, evaluating a continuous search space by rendering multi-camera configurations in a loop remains computationally unfeasible \cite{jorgensen2018simulation}. To resolve this computational bottleneck the continuous search space can be discretised by pre-rendering a dense network of $N$ candidate cameras surrounding the subject. This confines rendering to an upfront computational cost that is independent of the number of configurations evaluated. The discretisation of the search space makes this a combinatorial subset selection problem. The objective shifts to finding the optimal fixed-size subset of $k$ cameras from the broader set of $N$ cameras.
Even though the search space is restricted the number of possible combinations remains large. Evaluating countless multi-camera configurations requires a formalised optimisation approach. This formulation motivates the use of a discrete computational intelligence algorithms capable of maintaining a fixed-subset selection of cameras. 

\subsection{Problem Statement}
\label{sec:problemstatement}

Present literature on evaluating camera placement and geometry is limited to real world setups, marker-based systems or only considering a single pose estimation algorithm. This severely limits the exploration and evaluation of optimal camera configurations whilst ignoring the potential effects that the choice of pose estimation models have on a multiview markerless pose estimation system.

Markerless pose estimation offers solutions to the drawbacks of marker-based systems however markerless pose estimation systems have not seen widespread adoption. Further research is required to improve markerless pose estimation systems to further bridge the gap between gold standard marker-based systems and markerless pose estimation systems. Two major components that influence the performance of multiview markerless pose estimation systems are the human pose estimation algorithm and optical cameras. The choice of pose estimation algorithm can affect the system as different models have different strengths and properties. Optical cameras that capture a subject's motion sequence can influence the performance as both the camera and its placement within the system influence the quality of the data being processed. Little is understood about optimal camera configurations and setups for markerless pose estimation systems or even how the choice of algorithm might influence the camera setup.
\subsection{Research Aims and Objectives}
\label{sec:researchaimsandobjectives}

In light of the problem statement this research aims to better characterise optimal camera configurations for general action motion capture, action specific layouts across seven exercises and seeing if the choice of pose estimation algorithm influences the optimal camera configurations. This can be broken down into the following objectives.

Create a suitable synthetic data generation pipeline for rendering out data.
This first objective can be broken down into the following sub objectives.
\begin{enumerate}
    \item Select a suitable DCC in which renders will take place for the output of data. The choice of digital content creation platform influences the programming language being used, the tools that are available such as add-ons or extensions and could influence render time, computational resource consumption. 
    \item Select a human body model to represent a human in the 3D world. To be able to accurately create synthetic data an intelligent body model that can represent the human shape and underlying biomechanical skeletal structure is needed.
    \item Obtain suitable motion data to drive the human body model serving as ground truth data. Finding real world or realistic motion capture data in the domain of strength and conditioning exercises is crucial to improve multiview markerless pose estimation systems for sports scientists and other exercise professionals
    \item Create a camera rig layout to capture and render out the animation sequence. Given the continuous search space we need to discretise the environment by setting up cameras that surround the subject with sufficient coverage. 
\end{enumerate}

The second objective is to implement an optimisation algorithm to characterise and evaluate camera configurations. This second objective can be broken down into the following sub objectives
\begin{enumerate}
    \item Implement an optimisation algorithm capable of solving a fixed set-sized combinatorial selection problem. Systems differ in the number of cameras that are available to them. We therefore need to evaluate optimal camera configurations for varying fixed set sizes, $K$, where $K$ is the number of cameras available to the system.    
    \item Define an objective function that the optimisation algorithm can use to evaluate a camera configurations using metrics such as the Mean Per Joint Position Error (MPJPE) along with triangulability and visibility metrics. This is the most important aspect of the optimisation loop as it needs to be well defined to both reward and penalise the algorithm ensuring that it does not cheat its way to a "good result".
\end{enumerate}

The third objective is to test how choice of algorithm and contextual factors influence the choice of camera configuration. This third objective can be broken down into the following sub objectives
\begin{enumerate}
    \item Conduct inference on the rendered output data using different pose estimation algorithms to serve as input to the optimisation algorithm. Since the inferred keypoints are being compared to the synthetic ground truth data and used in the optimisation algorithm we can compare how the best camera configurations found by the optimisation algorithm differ by pose estimation algorithm. 
    \item Run the optimisation algorithm on subsets of the rendered data. Given that some specialists might only be interested in a particular class or exercise it would be beneficial to know how camera configurations differ when considering different exercises, as an example comparing pushups to jumping jacks.
\end{enumerate}